\chapter{数列}
\label{chap:sequences}

\begin{learninggoals}
  \begin{enumerate}
    \item 在通项、递推关系与前 $n$ 项和之间转换，并据此判断数列的单调性；
    \item 掌握等差、等比数列的结构与求和公式，理解公式成立的条件；
    \item 能从项的形状选择分组、错位相减或裂项相消；
    \item 能通过平移、倒数、除以指数项或特征根，把递推关系化为熟悉的数列；
    \item 会用可相消的放缩处理数列不等式，而不只停留在逐项估计。
  \end{enumerate}
\end{learninggoals}

\section{数列的语言与基本关系}

\begin{definitionbox}[数列、通项与递推关系]
  按照一定次序排列的一列数称为\textbf{数列}，记作
  \[
    a_1,a_2,\ldots,a_n,\ldots,
    \qquad\text{或}\qquad \{a_n\}.
  \]
  $a_n$ 是第 $n$ 项。若 $a_n$ 可由 $n$ 的式子直接表示，这个式子称为通项公式；若从已知的前若干项能够依次确定后项，这种关系称为递推关系。
\end{definitionbox}

数列可看作定义域为正整数集（或其有限子集）的函数。它既可以列表，也可以用通项、递推关系或点列 $(n,a_n)$ 表示。因为自变量只取整数，研究其增减时必须比较相邻两项，不能把连续函数的结论不加辨析地搬过来。

\begin{center}
\begin{tikzpicture}[x=.72cm,y=.34cm]
  \draw[-{Latex[length=2mm]},BookBlue] (0,0)--(8.3,0) node[right] {$n$};
  \draw[-{Latex[length=2mm]},BookBlue] (0,-2)--(0,5.2) node[above] {$a_n$};
  \foreach \x in {1,...,7} \draw[BookBlue!45] (\x,.11)--(\x,-.11) node[below=3pt,font=\scriptsize] {\x};
  \foreach \x in {1,...,7} {
    \pgfmathsetmacro{\y}{(\x-4.5)*(\x-4.5)/2-1.125}
    \fill[BookGold] (\x,\y) circle (2pt);
  }
  \draw[BookGold,dashed] plot[smooth,domain=.7:7.3,samples=60] (\x,{(\x-4.5)*(\x-4.5)/2-1.125});
  \node[align=left,font=\small,BookGreen!70!black] at (10,2.2) {虚线只帮助观察趋势；\\数列本身是离散的点。};
\end{tikzpicture}
\end{center}

\subsection{部分和与通项}

记
\[
  S_n=a_1+a_2+\cdots+a_n.
\]

\begin{lawbox}[由前 $n$ 项和恢复通项]
  \[
    a_1=S_1,\qquad a_n=S_n-S_{n-1}\quad(n\ge 2).
  \]
  第二个式子不能直接用于 $n=1$，除非另行约定 $S_0=0$。
\end{lawbox}

\begin{corollarybox}[判断单调性的两个入口]
  对相邻两项作差：
  \[
    a_{n+1}-a_n\gtreqless 0;
  \]
  当各项同号且分式、指数结构明显时，也可比较
  \[
    \frac{a_{n+1}}{a_n}\gtreqless 1.
  \]
  比值法必须先确认分母和各项的符号；符号不定时，作差通常更稳妥。
\end{corollarybox}

\begin{examplebox}[原稿型：离散自变量下的增减转折]
  设 $a_n=n^2-9n$，判断数列的单调性。

  \textbf{解：}比较相邻两项：
  \[
    a_{n+1}-a_n=(n+1)^2-9(n+1)-(n^2-9n)=2n-8.
  \]
  当 $1\le n<4$ 时差为负；$n=4$ 时差为零；$n>4$ 时差为正。因此
  \[
    a_1>a_2>a_3>a_4=a_5<a_6<a_7<\cdots.
  \]
  \textbf{辨析：}把 $a_n$ 配方为 $(n-\tfrac92)^2-\tfrac{81}{4}$ 可以看出趋势，但顶点横坐标不是整数；只有比较 $a_4,a_5$，才能准确写出离散数列的最低项。
\end{examplebox}

\section{等差数列}

\begin{definitionbox}[等差数列]
  从第 $2$ 项起，每一项与前一项的差等于同一个常数 $d$，即
  \[
    a_n-a_{n-1}=d\quad(n\ge2),
  \]
  则称 $\{a_n\}$ 为等差数列，$d$ 为公差。
\end{definitionbox}

把从 $a_1$ 到 $a_n$ 的相邻差相加，左端逐项抵消：
\[
  (a_2-a_1)+(a_3-a_2)+\cdots+(a_n-a_{n-1})=(n-1)d.
\]
由此得到通项公式。

\begin{lawbox}[等差数列的通项与前 $n$ 项和]
  \[
    a_n=a_1+(n-1)d,
    \qquad
    S_n=\frac{n(a_1+a_n)}2
       =na_1+\frac{n(n-1)}2d.
  \]
  第二个公式来自倒序相加：$S_n$ 与倒写后的 $S_n$ 对应相加，每一列都等于 $a_1+a_n$。
\end{lawbox}

\begin{corollarybox}[等差数列的对称性]
  若下标满足 $i+j=p+q$，则
  \[
    a_i+a_j=a_p+a_q.
  \]
  特别地，若 $i+j=2m$，则 $a_i+a_j=2a_m$。三个数 $a,b,c$ 成等差数列的充要条件是 $2b=a+c$。
\end{corollarybox}

\begin{corollarybox}[从任意一项表示和]
  对任意固定的 $m$，有
  \[
    S_n=na_m+\frac{n(n+1-2m)}2d.
  \]
  当 $m=(n+1)/2$ 时，第二项消失，这就是“项数乘中间项”；项数为偶数时，则使用中间两项的平均数。
\end{corollarybox}

\subsection{部分和的离散抛物线结构}

\[
  S_n=\frac d2n^2+\left(a_1-\frac d2\right)n.
\]
当 $d\ne0$ 时，$S_n$ 是关于 $n$ 的二次式，但只在正整数点上取值。求最大、最小值时，可以先找抛物线对称轴，再比较轴两侧相邻的整数点；也可以直接观察 $S_n-S_{n-1}=a_n$ 的符号变化。

\begin{corollarybox}[等差数列前 $n$ 项和的极值]
  \begin{itemize}
    \item 若 $d<0$，且 $a_m\ge0$、$a_{m+1}\le0$，则 $S_m$ 为最大值；
    \item 若 $d>0$，且 $a_m\le0$、$a_{m+1}\ge0$，则 $S_m$ 为最小值；
    \item 若某一项等于 $0$，相邻两个部分和可能同时取得极值。
  \end{itemize}
  本质是 $S_n$ 的相邻差恰为 $a_n$。
\end{corollarybox}

\begin{examplebox}[原稿例：平方条件揭示相邻项]
  等差数列 $\{a_n\}$ 的公差 $d\ne0$，前 $n$ 项和为 $S_n$，且
  \[
    a_2^2+a_3^2=a_4^2+a_5^2,\qquad S_7=7.
  \]
  求 $a_n$ 与 $S_n$。

  \textbf{解：}设公差为 $d$。移项并作平方差分解：
  \begin{align*}
    a_2^2-a_4^2&=a_5^2-a_3^2,\\
    (a_2-a_4)(a_2+a_4)&=(a_5-a_3)(a_5+a_3).
  \end{align*}
  因为 $a_2-a_4=-2d$、$a_5-a_3=2d$，且 $d\ne0$，约去 $2d$ 后得
  \[
    -(a_2+a_4)=a_5+a_3.
  \]
  利用下标对称，$a_2+a_4=2a_3$、$a_3+a_5=2a_4$，上式等价于
  \[
    a_3=-a_4.
  \]
  又 $S_7=7a_4=7$，故 $a_4=1$、$a_3=-1$，从而 $d=2$。因此
  \[
    \boxed{a_n=2n-7},\qquad
    \boxed{S_n=n^2-6n}.
  \]
  \textbf{方法启示：}平方条件先用平方差拆开，再把“项的和”转成下标和相同的对称关系；这比一开始完全展开 $a_1,d$ 更能暴露结构。
\end{examplebox}

\begin{examplebox}[原稿方法：由两组部分和之比求中间项之比]
  两个等差数列 $\{a_n\}$、$\{b_n\}$ 的前 $n$ 项和分别为 $S_n,T_n$。已知
  \[
    \frac{S_n}{T_n}=\frac{7n+1}{4n+27},
  \]
  求 $a_{11}/b_{11}$。

  \textbf{解：}等差数列连续奇数项的和等于“项数乘中间项”，所以
  \[
    S_{21}=21a_{11},\qquad T_{21}=21b_{11}.
  \]
  直接把 $n=21$ 代入已知比式：
  \[
    \frac{a_{11}}{b_{11}}
    =\frac{S_{21}}{T_{21}}
    =\frac{7\cdot21+1}{4\cdot21+27}
    =\boxed{\frac43}.
  \]
  \textbf{关键转化：}要求第 $m$ 项之比，应想到取前 $2m-1$ 项，使 $a_m,b_m$ 同时成为各自的中间项。
\end{examplebox}

\section{等比数列}

\begin{definitionbox}[等比数列]
  从第 $2$ 项起，每一项与前一项的比等于同一个非零常数 $q$：
  \[
    \frac{a_n}{a_{n-1}}=q\quad(n\ge2),
  \]
  则称 $\{a_n\}$ 为等比数列，$q$ 为公比。按此定义，等比数列的每一项都不为零。
\end{definitionbox}

\begin{lawbox}[等比数列的通项与前 $n$ 项和]
  \[
    a_n=a_1q^{n-1},
  \]
  并且
  \[
    S_n=
    \begin{cases}
      \dfrac{a_1(1-q^n)}{1-q},&q\ne1,\\[6pt]
      na_1,&q=1.
    \end{cases}
  \]
  当 $q\ne1$ 时，将 $S_n$ 与 $qS_n$ 错开相减，除首尾两项外全部抵消。$q=1$ 时不能除以 $1-q$，必须单独讨论。
\end{lawbox}

\begin{corollarybox}[等比数列的下标结构]
  \[
    a_n=a_mq^{n-m}.
  \]
  若 $i+j=p+q$，则 $a_i a_j=a_p a_q$；特别地，$i+j=2m$ 时有 $a_i a_j=a_m^2$。

  三个非零数 $a,b,c$ 成等比数列的充要条件是 $b^2=ac$ 且相邻两项比值相同。只写 $b^2=ac$ 时必须防止 $0,0,1$ 这类反例。
\end{corollarybox}

\begin{corollarybox}[等比数列的单调性]
  由
  \[
    a_{n+1}-a_n=a_1q^{n-1}(q-1)
  \]
  可知：当 $q>0$ 时结合 $a_1$ 与 $q-1$ 的符号判断增减；$q=1$ 时为常数列；$q<0$ 时各项正负交替，通常不是单调数列。不能只看 $q$ 是否大于 $1$。
\end{corollarybox}

\begin{corollarybox}[等比分组仍保留等比结构]
  若 $\{a_n\}$ 为等比数列，把相邻的 $m$ 项分为一组，则各组的和
  \[
    S_m,\quad S_{2m}-S_m,\quad S_{3m}-S_{2m},\ldots
  \]
  依次相差因子 $q^m$，因而在各组非零时构成等比数列。
\end{corollarybox}

\begin{examplebox}[原稿例：带线性系数的等比求和]
  求
  \[
    S_n=\frac12+2\left(\frac12\right)^2
       +3\left(\frac12\right)^3+\cdots
       +n\left(\frac12\right)^n.
  \]

  \textbf{解：}乘以公比 $1/2$，使同次幂错位：
  \begin{align*}
    S_n&=\frac12+2\left(\frac12\right)^2+cdots+n\left(\frac12\right)^n,\\
    \frac12S_n&=\left(\frac12\right)^2+2\left(\frac12\right)^3+cdots
      +(n-1)\left(\frac12\right)^n+n\left(\frac12\right)^{n+1}.
  \end{align*}
  相减后，中间各项的系数都变成 $1$：
  \begin{align*}
    \frac12S_n
      &=\frac12+\left(\frac12\right)^2+\cdots+\left(\frac12\right)^n
        -n\left(\frac12\right)^{n+1}\\
      &=1-\frac1{2^n}-\frac n{2^{n+1}}.
  \end{align*}
  所以
  \[
    \boxed{S_n=2-\frac{n+2}{2^n}}.
  \]
  \textbf{方法启示：}“等差系数乘等比项”应先乘公比再错位相减；相减的目的，是把线性系数降为常数。
\end{examplebox}

\section{数列求和的结构}

\subsection{常用有限和}

\begin{lawbox}[三组基本公式]
  \begin{align*}
    1+2+\cdots+n&=\frac{n(n+1)}2,\\
    1^2+2^2+\cdots+n^2&=\frac{n(n+1)(2n+1)}6,\\
    1^3+2^3+\cdots+n^3&=\left[\frac{n(n+1)}2\right]^2.
  \end{align*}
\end{lawbox}

平方和公式可以从
\[
  (k+1)^3-k^3=3k^2+3k+1
\]
对 $k=1,2,\ldots,n$ 求和得到；立方和公式则可从四次方差分出发。这里的共同思想是：把目标项嵌进一个能首尾相消的差分式。

\subsection{裂项相消}

\begin{corollarybox}[常见裂项形状]
  \begin{align*}
    \frac1{n(n+1)}&=\frac1n-\frac1{n+1},\\
    \frac1{n(n+k)}&=\frac1k\left(\frac1n-\frac1{n+k}\right),\\
    \frac1{(n-1)n(n+1)}&=\frac12\left(\frac1{(n-1)n}-\frac1{n(n+1)}\right),\\
    \frac1{\sqrt{n+1}+\sqrt n}&=\sqrt{n+1}-\sqrt n,\\
    \frac{n}{(n+1)!}&=\frac1{n!}-\frac1{(n+1)!}.
  \end{align*}
  裂项不是机械背公式：先观察分母相邻因子的差，再让拆开的两项在连续求和中对齐。
\end{corollarybox}

\begin{examplebox}[原稿例：递推关系直接制造裂项]
  数列 $\{a_n\}$ 满足 $a_1=1$，
  \[
    a_{n+1}=a_n^2+a_n.
  \]
  求和
  \[
    \sum_{k=1}^{n}\frac1{a_k+1}.
  \]

  \textbf{解：}递推式可写成 $a_{n+1}=a_n(a_n+1)$，因此
  \[
    \frac1{a_n+1}
    =\frac{a_n}{a_n(a_n+1)}
    =\frac1{a_n}-\frac1{a_{n+1}}.
  \]
  连续相加后只留下首尾：
  \[
    \boxed{\sum_{k=1}^{n}\frac1{a_k+1}=1-\frac1{a_{n+1}}}.
  \]
  \textbf{关键转化：}题目没有给出 $a_n$ 的通项，也不需要先求通项；递推式已经提供了相邻两项之间的裂项接口。
\end{examplebox}

\section{递推数列的构造}

递推式的目标不是一律“展开到第一项”，而是先识别不变的结构：差、比、平移后的比，或两个不动点之间的比例。

\begin{corollarybox}[累加与累乘]
  若 $a_{n+1}-a_n=f(n)$，则
  \[
    a_n=a_1+\sum_{k=1}^{n-1}f(k).
  \]
  若 $a_{n+1}/a_n=g(n)$ 且各分母非零，则
  \[
    a_n=a_1\prod_{k=1}^{n-1}g(k).
  \]
  前者寻找相邻差，后者寻找相邻比；连加或连乘时要同时检查起止下标。
\end{corollarybox}

\subsection{一阶线性递推：平移或除以指数项}

\begin{lawbox}[常数型一阶递推]
  对
  \[
    a_{n+1}=pa_n+c\qquad(p\ne1),
  \]
  取不动点 $L=c/(1-p)$，则
  \[
    a_{n+1}-L=p(a_n-L).
  \]
  因而 $\{a_n-L\}$ 是等比数列。所谓“配常数”，本质是把递推式的固定点移到原点。
\end{lawbox}

当非齐次项是 $n$ 的多项式时，可尝试构造 $a_n+An+B$ 或更高次多项式；当非齐次项是 $r^n$ 时，可尝试构造 $a_n+Ar^n$，也可直接除以合适的指数项。若 $r=p$ 出现“共振”，通常要把试探项提高为 $Anr^n$。

\begin{examplebox}[原稿型：加入一次式构造等比数列]
  已知 $a_1=2$，$a_{n+1}=3a_n+2n$。求通项公式。

  \textbf{解：}设 $b_n=a_n+An+B$，希望 $b_{n+1}=3b_n$。比较
  \begin{align*}
    b_{n+1}&=3a_n+2n+A(n+1)+B,\\
    3b_n&=3a_n+3An+3B,
  \end{align*}
  得
  \[
    2+A=3A,qquad A+B=3B,
  \]
  所以 $A=1$、$B=\tfrac12$。于是
  \[
    a_{n+1}+n+\frac32=3\left(a_n+n+\frac12\right),
  \]
  即 $\{a_n+n+\tfrac12\}$ 是首项 $7/2$、公比 $3$ 的等比数列。因此
  \[
    \boxed{a_n=\frac72\,3^{n-1}-n-\frac12}.
  \]
  \textbf{方法启示：}待定的一次式不是凭经验硬凑；把目标关系写出后比较 $n$ 的系数与常数项，就能系统确定。
\end{examplebox}

\subsection{分式递推：围绕不动点变换}

\begin{lawbox}[分式递推的不动点法]
  对 $a_{n+1}=f(a_n)$，先解 $f(x)=x$。若有两个不同不动点 $\alpha,\beta$，常可整理成
  \[
    \frac{a_{n+1}-\alpha}{a_{n+1}-\beta}
      =q\frac{a_n-\alpha}{a_n-\beta},
  \]
  从而把原递推化成等比数列。若其中一个不动点使式子更简单，也可只取倒数 $1/(a_n-\alpha)$ 再平移。
\end{lawbox}

\begin{examplebox}[原稿例：分式递推的倒数平移]
  已知 $a_1=2$，
  \[
    a_{n+1}=\frac{2a_n+2}{a_n+3}.
  \]
  求 $a_n$。

  \textbf{解：}方程 $(2x+2)/(x+3)=x$ 的两个不动点为 $1,-2$。先围绕 $1$ 整理：
  \[
    a_{n+1}-1=\frac{a_n-1}{a_n+3}.
  \]
  取倒数，并令 $b_n=1/(a_n-1)$，则
  \[
    b_{n+1}=\frac{a_n+3}{a_n-1}=1+4b_n.
  \]
  再平移一次：
  \[
    b_{n+1}+\frac13=4\left(b_n+\frac13\right).
  \]
  因为 $b_1=1$，所以
  \[
    b_n+\frac13=\frac43\,4^{n-1}=\frac{4^n}{3}.
  \]
  故
  \[
    \boxed{a_n=1+\frac{3}{4^n-1}}.
  \]
  \textbf{方法启示：}第一次变换利用不动点消去常数，第二次平移把线性递推化为等比；两步分别解决不同的障碍。
\end{examplebox}

\begin{examplebox}[原稿例：自加权递推先比较相邻两式]
  设 $a_1=1$，且对 $n\ge2$，
  \[
    a_n=a_1+\frac12a_2+\frac13a_3+\cdots+\frac1{n-1}a_{n-1}.
  \]
  求 $a_n$，并确定满足 $a_n>2007$ 的最小正整数 $n$。

  \textbf{解：}把 $n$ 与 $n+1$ 对应的两式相减：
  \[
    a_{n+1}-a_n=\frac1n a_n,
  \]
  即
  \[
    \frac{a_{n+1}}{a_n}=\frac{n+1}{n}.
  \]
  又由原式知 $a_2=a_1=1$，连续相乘得
  \[
    \frac{a_n}{a_2}=\frac32\cdot\frac43\cdots\frac n{n-1}=\frac n2,
  \]
  所以 $a_1=1$，且 $n\ge2$ 时 $a_n=n/2$。因此 $a_n>2007$ 等价于 $n>4014$，最小的 $n$ 为
  \[
    \boxed{4015}.
  \]
  \textbf{方法启示：}含有逐步延长的和式时，先写相邻两个下标的等式并相减；长和式会只剩新添的一项。
\end{examplebox}

\subsection{二阶线性递推与特征根}

\begin{lawbox}[二阶常系数递推]
  对
  \[
    a_{n+2}=s a_{n+1}+t a_n,
  \]
  设法构造
  \[
    a_{n+1}+\lambda a_n=(s+\lambda)(a_n+\lambda a_{n-1}).
  \]
  比较系数得
  \[
    \lambda^2+s\lambda-t=0.
  \]
  这与设试探解 $a_n=r^n$ 得到的特征方程 $r^2-sr-t=0$ 等价（$r=s+\lambda$）。若有两个不同实根 $r_1,r_2$，则
  \[
    a_n=C_1r_1^{n-1}+C_2r_2^{n-1},
  \]
  常数由初值确定。
\end{lawbox}

\begin{examplebox}[原稿方法：同时构造两条一阶递推]
  对递推关系
  \[
    a_n=2a_{n-1}+3a_{n-2}\qquad(n\ge3),
  \]
  说明如何降阶。

  \textbf{解：}特征方程
  \[
    r^2-2r-3=0
  \]
  的根为 $3,-1$。不直接写结论，而分别把原递推整理为
  \begin{align*}
    a_n+a_{n-1}&=3(a_{n-1}+a_{n-2}),\\
    a_n-3a_{n-1}&=-(a_{n-1}-3a_{n-2}).
  \end{align*}
  因此 $\{a_n+a_{n-1}\}$ 与 $\{a_n-3a_{n-1}\}$ 分别是公比为 $3$ 与 $-1$ 的等比数列。代入题目给定的两个初值后，两式联立即可求出 $a_n$。

  \textbf{方法启示：}两个特征根对应两条独立的一阶关系。保留这一步构造，可以看清通项中两个指数项从何而来，而不是把特征根公式当作黑箱。
\end{examplebox}

\section{放缩与数列不等式}

放缩的目的不是把每一项估得尽可能接近，而是把它改造成能求和、最好能相消的形式。先确定所需方向，再寻找相邻下标之间的差。

\begin{corollarybox}[设计可相消放缩的步骤]
  \begin{enumerate}
    \item 先写清需要的是上界还是下界；
    \item 猜测相邻差 $F(k+1)-F(k)$，使求和后只留下边界；
    \item 检查不等号方向与适用的起始下标；
    \item 必要时把前一两项单独处理，再对尾项统一放缩。
  \end{enumerate}
\end{corollarybox}

\begin{examplebox}[原稿例：用裂项证明平方倒数和有界]
  证明
  \[
    1+\frac1{2^2}+\frac1{3^2}+\cdots+\frac1{n^2}<2
    \qquad(n\ge2).
  \]

  \textbf{证明：}对 $k\ge2$，有 $k^2>k(k-1)$，故
  \[
    \frac1{k^2}<\frac1{k(k-1)}=\frac1{k-1}-\frac1k.
  \]
  所以
  \[
    1+\sum_{k=2}^{n}\frac1{k^2}
    <1+\sum_{k=2}^{n}\left(\frac1{k-1}-\frac1k\right)
    =2-\frac1n<2.
  \]
  \textbf{方法启示：}逐项上界 $1/[k(k-1)]$ 虽不精细，却能完全相消；整体结构比单项误差更重要。
\end{examplebox}

\begin{examplebox}[原稿例：平方和尾项的定向放缩]
  证明
  \[
    \frac1{2^2}+\frac1{3^2}+\cdots+\frac1{n^2}<\frac34
    \qquad(n\ge2).
  \]

  \textbf{证明：}首项单独保留。对 $k\ge3$，因为 $k^2>k(k-1)$，
  \[
    \frac1{k^2}<\frac1{k(k-1)}=\frac1{k-1}-\frac1k.
  \]
  因此
  \[
    \sum_{k=2}^{n}\frac1{k^2}
    <\frac14+\sum_{k=3}^{n}\left(\frac1{k-1}-\frac1k\right)
    =\frac34-\frac1n<\frac34.
  \]
  \textbf{辨析：}若从 $k=2$ 就统一采用同一上界，只能得到较弱的常数。把少数边界项单独处理，常能显著改善整体估计。
\end{examplebox}

\begin{examplebox}[原稿例：根式项的两种相邻放缩]
  设
  \[
    T_n=1+\frac1{\sqrt2}+\frac1{\sqrt3}+\cdots+\frac1{\sqrt n}.
  \]
  证明
  \[
    2(\sqrt{n+1}-1)<T_n<2\sqrt n.
  \]

  \textbf{证明：}对 $k\ge1$，有
  \[
    \frac1{\sqrt k}
      =\frac2{2\sqrt k}
      >\frac2{\sqrt{k+1}+\sqrt k}
      =2(\sqrt{k+1}-\sqrt k).
  \]
  求和即得 $T_n>2(\sqrt{n+1}-1)$。

  对 $k\ge2$，又有
  \[
    \frac1{\sqrt k}
      <\frac2{\sqrt k+\sqrt{k-1}}
      =2(\sqrt k-\sqrt{k-1}).
  \]
  把首项 $1$ 单独保留并对其余项求和：
  \[
    T_n<1+2(\sqrt n-1)=2\sqrt n-1<2\sqrt n.
  \]
  \textbf{方法启示：}同一个根式项，上界和下界要分别选择前后相邻的根式差；放缩方向改变，配对的下标也随之改变。
\end{examplebox}
